\begin{table}[!htbp]\centering\small
\caption{Rents of insiders and entrants against incomes}\label{tab:facts}
\begin{adjustbox}{max width=\linewidth}\begin{tabular}{lllcccc}\toprule
Rent measure & Period & Units & Rent growth & Income growth & Difference & Share with rent $>$ income\\
 & & & \multicolumn{3}{c}{(log points, weighted)} & \\\midrule
All leases (IPVA), Spain & 2015--2023 & 703 municipalities & 1650.0 & 31.0 & 1619.0 & 0\,\%\\
New leases (deposits), Catalonia & 2015--2023 & 289 municipalities & 39.0 & 27.3 & 11.7 & 77\,\%\\
New leases per m$^2$ (deposits), Basque Country & 2016--2023 & 67 municipalities & 21.2 & 26.4 & $-$5.2 & 10\,\%\\
Deposits on new leases, Comunitat Valenciana & 2020--2023 & 104 municipalities & 22.9 & 15.9 & 7.0 & --\\
\addlinespace
\multicolumn{7}{l}{\textit{Entry gap: new leases relative to leases in force}}\\
INE index, new vs existing, Spain & 2021--2024 & national & 17.6 & 7.2 & 10.4 & \\
AEAT, new vs all leases, per euro of cadastral value & 2024 & 402 municipalities & & & 15.9 & 99.5\,\%$^a$\\
\bottomrule\end{tabular}\end{adjustbox}
\begin{minipage}{0.97\linewidth}\footnotesize\vspace{4pt}
\textit{Notes}: income is the median net income per consumption unit of the municipality (INE household income atlas). Means weighted by 2015 population. IPVA: index of all leases in force (tax returns). Deposit registers record leases at signature: mean rent in Catalonia (contracts of more than one year), rent per m$^2$ built in the Basque Country (free-market collective main residences), deposit amount in the Comunitat Valenciana (one month of rent by law; partial coverage). AEAT: tax statistics of dwellings declared in 2024, municipalities above 20,000 inhabitants; the entry gap is $\ln(\text{rent}_{new}/\text{value}_{new})-\ln(\text{rent}_{all}/\text{value}_{all})$. $^a$ Share of municipalities with a positive gap.
\end{minipage}\end{table}